\begin{singlespace}
\small
\enlargethispage{2cm} % Sécurité absolue pour forcer tout le bloc sur une seule page
\setlength{\parskip}{1.5pt}

\phantomsection
\addcontentsline{toc}{section}{Résumé}
\begin{center}
    \Large\bfseries Résumé
\end{center}
\vspace{0.1cm}

Ce projet de fin d'études a été réalisé au sein de l'entreprise BIOSPHERE dans le cadre de la mise en place d'un système décisionnel intelligent dédié à l'analyse et à la prévision des stocks.

L'objectif principal consiste à transformer les données opérationnelles de l'entreprise en informations stratégiques permettant d'améliorer la prise de décision, d'anticiper les risques de rupture de stock et d'optimiser les opérations de réapprovisionnement.

Pour atteindre cet objectif, une architecture décisionnelle complète a été mise en œuvre. Celle-ci repose sur la création d'un Data Warehouse alimenté par un processus ETL automatisé développé à l'aide de n8n, Docker et SQL Server. Les données intégrées ont ensuite été exploitées dans un environnement de Business Intelligence afin de produire des indicateurs de performance et des tableaux de bord décisionnels.

Parallèlement, deux axes de Data Science ont été développés. Le premier repose sur des modèles de classification permettant d'identifier les produits présentant un risque de rupture de stock. Le second s'appuie sur des modèles de régression destinés à prévoir les niveaux futurs de stock et à générer des recommandations de réapprovisionnement.

Enfin, l'ensemble des composants décisionnels et prédictifs a été intégré au sein d'une application desktop développée en Python, offrant aux utilisateurs une interface unifiée regroupant les tableaux de bord, les prévisions de stock, les alertes intelligentes et les fonctionnalités d'administration.

Les résultats obtenus démontrent l'intérêt de l'intégration des techniques de Business Intelligence et de Machine Learning pour renforcer les capacités d'analyse et de pilotage des activités de l'entreprise BIOSPHERE.

\vspace{0.2cm}
\noindent\textbf{Mots-clés :} Business Intelligence, Data Warehouse, ETL, Machine Learning, Classification, Régression, Prévision des stocks, Aide à la décision, Application Desktop.

\vspace{0.5cm} 

\phantomsection
\addcontentsline{toc}{section}{Abstract}
\begin{center}
    \Large\bfseries Abstract
\end{center}
\vspace{0.1cm}

This final-year project was carried out within BIOSPHERE Company and focuses on the development of an intelligent decision support system dedicated to inventory analysis and stock forecasting.

The main objective is to transform operational data into strategic information that supports decision-making, anticipates stock-out risks, and improves replenishment planning.

To achieve this objective, a complete decision-making architecture was implemented. It relies on a Data Warehouse supplied through an automated ETL process developed using n8n, Docker, and SQL Server. The integrated data were then exploited through Business Intelligence tools to generate performance indicators and decision dashboards.

In parallel, two Data Science components were developed. The first component is based on classification models designed to identify products at risk of stock shortage. The second component relies on regression models used to forecast future stock levels and generate replenishment recommendations.

Finally, all analytical and predictive components were integrated into a desktop application developed in Python, providing users with a unified interface that includes dashboards, stock forecasts, intelligent alerts, and administration functionalities.

The obtained results demonstrate the value of combining Business Intelligence and Machine Learning techniques to enhance operational monitoring and decision-making capabilities within BIOSPHERE.

\vspace{0.2cm}
\noindent\textbf{Keywords:} Business Intelligence, Data Warehouse, ETL, Machine Learning, Classification, Regression, Stock Forecasting, Decision Support System, Desktop Application.

\end{singlespace}